\chapter{AI Assistance and Retrospective Prompt Summaries}
\label{app:prompt-record}

\section{Implementation and Experimentation}
I used AI assistance to review existing code and results, develop implementations,
troubleshoot failures and prepare comparisons of individual models and fusion methods.
The tasks included checking label meaning, modality correspondence, grouped partitions,
training-side fitting of fusion, fixed experiment settings and the preservation of
configurations, seeds, predictions and failed attempts. The prompts also concerned
checking required gates before execution, evaluating planned missing or degraded inputs,
keeping comparisons within the same study protocol, and retaining unsuccessful attempts
without changing settings to improve a held-out score. These are retrospective task
summaries; the configurations and run records establish which experiments were completed.

\section{Retrieving and Checking Results}
AI assistance helped retrieve saved metrics and figures, relate them to their
experiments, check sample identities and compare reported numbers with available
predictions. It also supported descriptions of the models, negative or inconclusive
findings, incomplete work and missing evidence. These checks used the records available
for each revision and did not establish independent reproduction of model training.

\section{Archiving and Version Control of Experimental Outputs}
\label{sec:results-archiving}
AI assistance also supported planning for local archiving and Git version control of
experimental outputs. I retained these as planning tasks; the records do not establish completed
transfer or upload.

\section{Provenance}
The summaries describe tasks documented in retained requests and revision records.
Complete original implementation chats are unavailable, so the exact prompt wording
and sequence cannot be reconstructed. Earlier model versions were not recorded.
I therefore present task descriptions rather than attributing reconstructed dialogue
to an original exchange.

\section{Targeted Scientific Revision}
Codex assisted with strengthening the architectural justification, comparing the
findings with the literature, clarifying dataset-specific interpretations, and retaining
effect sizes, uncertainty and relevant investigation history. The revision used existing
evidence and preserved the distinction between completed, inconclusive and unexecuted
studies. It did not involve new model training.

\section{Evidence Checks and Revision of the Empirical Chapters}
I used the supplied chapter review to guide Codex-assisted checks of the empirical chapters
against available saved results and implementation records. The work covered study
roles, quantitative results, chapter structure, evidence tables, cross-references and
duplication. Codex assisted with revising the LaTeX text and checking saved evidence
without new model training or inference. The resulting account distinguishes directly
inspected records, findings available only in retained reports, and evidence that
remained unavailable.
